\section{Nuevos vectores de ataque que traen los AI Agent}

\subsection{Prompt Injection}

Ocultar o inyectar instrucciones engañosas en un sistema de IA para desviarlo de su comportamiento previsto. Esta es la amenaza de seguridad más central que enfrentan los AI Agent.

\subsection{Tool Misuse}

Manipular a un Agent mediante prompts engañosos para que abuse de las herramientas que tiene integradas. Por ejemplo, inducir al Agent a ejecutar operaciones de archivos o llamadas a API no autorizadas.

\subsection{Intent Breaking}

Manipular el plan de ejecución (plan) del Agent para redirigir sus acciones de la intención original hacia el objetivo del atacante.

\subsection{Identity Spoofing}

Aprovechar mecanismos de autenticación comprometidos para hacerse pasar por un Agent legítimo.

\subsection{Code Attacks}

Aprovechar la capacidad del Agent de ejecutar código para obtener acceso no autorizado al entorno de ejecución.

\begin{warningbox}{La superficie de ataque de los AI Agent es mucho mayor que la del software tradicional}
La superficie de ataque del software tradicional se concentra principalmente en la validación de entradas y la autenticación. La de los AI Agent, en cambio, se extiende a múltiples dimensiones: el prompt (que puede ser objeto de inyección), las llamadas a herramientas (que pueden ser objeto de abuso), el plan de ejecución (que puede ser manipulado), la autenticación de identidad (que puede ser suplantada) y el entorno de ejecución de código (que puede ser explotado).
\end{warningbox}

\subsection{Resumen de la sección}

Los AI Agent introducen al menos cinco nuevos tipos de vectores de ataque. La protección requiere actuar simultáneamente en múltiples capas: filtrado de prompts, control de permisos de herramientas, validación de planes, autenticación de identidad y aislamiento mediante sandboxing.

